% AUTO-GENERATED by scripts/generate_manuscript_figures.py -- do not edit by hand.
% Each caption is derived from the plotted numbers.

% ---- fig1: Activation frontier (figures/manuscript_v2/fig1_activation_frontier.pdf)
\caption{\textbf{Activation frontier under open exit.} Agent-based sweep at $\delta_0=0$ (exit open), $n=1890$ runs (9 $\sigma$ $\times$ 7 $\pi$ cells $\times$ 30 seeds per cell). (A) Mean maximum punish rate, the continuous enforcement-intensity surface. (B) Probability that a run is classified active, recomputed with the hierarchical classifier (active: max punish rate $\geq$ 0.1); overall P(active) $=$ 0.598. Enforcement activates only above a joint $\sigma$--$\pi$ threshold: the $\pi=0.01$ column (dashed outline) is dead in every one of its 9 cells, so an enforcement reward below that level cannot sustain an enforcement apparatus at any level of code observability.}

% ---- fig2: Exit closure is necessary for capture (figures/manuscript_v2/fig2_two_dimensions.pdf)
\caption{\textbf{Exit closure is necessary for capture.} All panels share one colour scale; 9 $\sigma$ $\times$ 7 $\pi$ cells $\times$ 30 seeds per cell. (A) Probability a run is active with the exit open ($\delta_0=0$, $n=1890$ runs): code geometry alone sets activation, overall P(active) $=$ 0.598. (B) Probability of capture in those same 1890 runs: \textbf{0 of 1890 runs reach capture with the exit open}, in every one of the 63 cells, despite that active rate of 0.598 -- an active enforcement apparatus never converts to capture while members can leave. (C) Probability of capture with the exit sealed ($\delta_0=0.95$, $n=1890$ runs): the same active region now captures. Panels B and C are the same quantity under the two exit conditions, so the contrast isolates the effect of closing the exit. The two dimensions are not fully orthogonal: circled cells in (C) are the 12 of 63 cells where sealing also raises P(active) (mean shift over all cells +0.062; +0.328 among circled cells), and activation falls in no cell. Regimes recomputed with the hierarchical classifier.}

% ---- fig3: Exit-closure dose-response (figures/manuscript_v2/fig3_exit_closure_dose_response.pdf)
\caption{\textbf{Exit-closure dose-response and the necessity of code geometry.} Capture rate against imposed exit closure $\delta_0$ ($n=1620$ runs total; 270 runs per pooled point, 30 per cell-specific point), with Wilson 95\% intervals. Pooled over all $\sigma\times\pi$ cells (circles) capture rises monotonically from 0.000 at $\delta_0=0$ to 0.889 at $\delta_0=0.95$. In the $\sigma=0.25$, $\pi=0.05$ cell (squares, dashed) the capture rate is 0.000 at every level including $\delta_0=0.95$: closing the exit cannot manufacture capture where the code geometry does not activate. Regimes recomputed with the hierarchical classifier.}

% ---- fig4: Privilege ablation (figures/manuscript_v2/fig4_privilege_ablation.pdf)
\caption{\textbf{Concentration is manufactured by the privilege architecture.} Role-independent concentration for the ablation floor (no privilege) and ceiling (full privilege bundle), 120 seeds per arm ($n=240$ runs); points are individual runs, black bars arm means. (A) Top-5\% share of punishment among active agents rises from 0.091 at the floor to 0.580 at the ceiling, so 84\% of ceiling concentration is privilege-manufactured. (B) Gini rises from 0.240 to 0.901 against random-allocation nulls (dotted) of 0.028 and 0.062; the floor sits close to its null, so concentration is not an emergent property of the unprivileged model. A committed random-allocation null exists only for the Gini statistic, so panel A carries no null reference.}

% ---- fig5: Circularity demonstration (figures/manuscript_v2/fig5_circularity_demonstration.pdf)
\caption{\textbf{Specification check for circular closure.} A specification in which outside-option closure is driven by enforcer share yields capture by construction, because that quantity mechanically drives retention, one of the two capture criteria ($\delta_{target} = \min(1, \delta_{baseline} + $ enforcer share$)$); the coupled arm ($n=4050$ runs) has a capture rate of 0.711. Exit capacity is therefore treated as an independent structural input. When closure is decoupled from enforcer share and driven by punishment intensity instead, capture does not occur at any coupling strength: 0 of 4050 runs at each of $k=1.5$, $3.0$ and $6.0$. Wilson 95\% intervals are shown; regimes are computed with the hierarchical classifier, under which capture is retained AND active and enforcer share is not itself a criterion.}
